% Part: history
% Chapter: biographies 
% Section: alonzo-church

\documentclass[../../../include/open-logic-section]{subfiles}

\begin{document}

\olfileid{his}{bio}{chu}

\olsection{الونزو چرچ}

\olphoto{church-alonzo}{الونزو چرچ}

الونزو چرچ د 1903 د جون پر 14مه په واشنګټن، ډي سي
کښې زېږېدلے و۔ په لومړني ماشومتوب کښې د يو هوايي
ټوپک پېښې د چرچ يوه سترګه ړنده کړه۔ هغه په 1920
کښې په کنېکټيکټ کښې د پوهنتون د چمتووالي ښوونځے
بشپړ کړ او په هماغه کال ئې په پرنسټن کښې د
پوهنتون زده کړه پيل کړه۔ دوکتورا ئې په 1927
کښې بشپړه کړه۔ له دوو کلونو بهر اوسېدو وروسته
چرچ پرنسټن ته راوګرځېد۔ چرچ د ډېر ادب او
احتياط له امله پېژندل کېدو۔ پر تخته ليکل ئې
بې‌عيبه وو، او مهمې مقالې به ئې په ډېر احتياط
په ډوکو سېمنټ، چې يو روڼ سريښ دے، پوښلې او
ساتلې۔ له علمي کارونو بهر ئې د ساينسي خيالي ادب
مجلې په شوق لوستلې، او که په ليکنو کښې ئې
ناسمې خبرې وليدې، مديرانو ته ئې په ليکلو
کښې هېڅ وېره نۀ کوله۔

د چرچ علمي لاسته راوړنې لويې وې۔ له خپلو
زده کوونکو سټيفن کليني او بارکلي روسر سره
ئې په ګډه د اغېزمنې محاسبې يوه تيوري،
لامبډا حساب، جوړه کړه، د الَن ټيورينګ له
خوا د ټيورينګ ماشين له جوړولو څخه په
خپلواک ډول۔ د محاسبه کېدنې دا دوه
تعريفونه معادل دي، او هغه څه ترې
راولاړېږي چې اوس \emph{د چرچ--ټيورينګ اصل}
نومېږي: د طبيعي عددونو يوه تابع په
اغېزمنه توګه محاسبه کېدونکې ده هله او
يوازې هله چې د ټيورينګ ماشين، يا
لامبډا حساب، په ذريعه محاسبه کېدونکې
وي۔ هغه دا هم ثابت کړل چې اوس
\emph{د چرچ قضيه} نومېږي: د لومړي
ترتيب د فارمولونو د اعتبار د پرېکړې
مسئله نۀ شي حل کېداے۔

چرچ تر زړښت پورې خپل کار جاري وساتۀ۔
په 1967 کښې ئې پرنسټن د UCLA دپاره
پرېښود، او هلته تر خپل تقاعد، په
1990 کښې، پروفېسر و۔ چرچ د 1995
د اګست پر 11مه د 92 کلونو په عمر
وفات شو۔

\begin{reading}
د چرچ د لنډ ژوندليک دپاره
\citet{EndertonND} وګورئ۔ د لامبډا
حساب او د پرېکړې مسئلې، يعنې
Entscheidungsproblem، په هکله د
چرچ اصلي ليکنې
\citet{Church1936,Church1936a} دي۔
\citet{Aspray1984} له چرچ سره يوه
مرکه ثبتوي چې په 1930s کښې د
پرنسټن د رياضياتو د ټولنې په
هکله ده۔ چرچ له 1936 څخه تر
1979 پورې د \emph{د سمبوليک
منطق مجلې} د کتابونو د کتنو
يوه لړۍ وليکله۔ دا ټولې د
جان مکفارلېن پر وېبپاڼه
خوندي شوې دي \citep{MacFarlane2015}۔
\end{reading}
\end{document}
